\section{Tests statistiques mobilisés et sources}
\label{ann:G}

Cette annexe récapitule les tests et procédures d'inférence utilisés dans le
mémoire, leur logique, la règle de décision et la source méthodologique de
référence. Sauf mention contraire, les erreurs-types sont clusterisées au
niveau de la grappe (unité primaire de sondage).

\subsection{Équilibre des covariables par la différence standardisée moyenne (SMD)}

Avant et après appariement, l'équilibre de chaque covariable $X$ entre
bénéficiaires ($D=1$) et témoins ($D=0$) est mesuré par la différence
standardisée moyenne
\[
  \mathrm{SMD}(X) = \frac{\bar{X}_{1} - \bar{X}_{0}}
       {\sqrt{\tfrac{1}{2}\left(s_{1}^{2} + s_{0}^{2}\right)}},
\]
où $\bar{X}_{d}$ et $s_{d}^{2}$ sont la moyenne et la variance dans le groupe
$d$. Un biais absolu inférieur à 5\,\% après appariement est le seuil
d'équilibre usuellement retenu. La statistique globale (biais moyen,
pseudo-$R^2$ du probit ré-estimé sur l'échantillon apparié) est produite par
la commande \texttt{pstest}. \textit{Sources :}
\textcite{rosenbaum1983,austin2009}.

\subsection{Estimation probit du score de propension}

Le score de propension $e(X)=\Pr(D=1\mid X)$ est estimé par un modèle probit
sur les covariables pré-traitement (taille et composition du ménage,
dépense par tête, milieu, région, sexe, âge et statut matrimonial du chef,
niveau d'instruction). La qualité d'ajustement est appréciée par le
pseudo-$R^2$ et le test du rapport de vraisemblance (LR $\chi^2$).
\textit{Source :} \textcite{rosenbaum1983}.

\subsection{Appariement sur le score de propension}

Trois algorithmes sont comparés~: plus proches voisins ($k=4$, avec remise),
noyau d'Epanechnikov et rayon (caliper, $\epsilon = 0{,}05$, sans
remise), mis en \oe uvre avec \texttt{psmatch2}. La zone de support commun
est imposée pour garantir la comparabilité. \textit{Sources :}
\textcite{rosenbaum1983,heckman1997,heckman1998,caliendo2008}.

\subsection{Estimateur de double différence}

L'effet du traitement sur les traités (ATT) est estimé par la spécification
d'interaction
\[
  Y_{it} = \alpha + \beta\, t + \gamma\, D_i + \delta\,(t \times D_i)
           + \varepsilon_{it},
\]
où $\delta$ est l'ATT. Sur l'échantillon apparié, la régression est pondérée
par les poids d'appariement (PSM-DD)~; le test de significativité de $\delta$
est un test de Student sur la combinaison linéaire $1.t\#1.D$
(\texttt{lincom}). \textit{Sources :} \textcite{card1994} pour la double
différence~; \textcite{heckman1997,heckman1998} pour sa combinaison avec
l'appariement.

\subsection{Erreurs-types clusterisées}

Toutes les inférences reposent sur des erreurs-types robustes à
l'hétéroscédasticité et à la corrélation intra-grappe
(\texttt{vce(cluster grappe)}), qui tiennent compte du plan de sondage à deux
degrés sans recourir aux poids de sondage. \textit{Source :}
\textcite{deaton1997}.

\subsection{Test placebo par permutation}

Pour vérifier l'absence d'effet artefactuel, un faux statut de traitement est
attribué aléatoirement à la même proportion de ménages non-bénéficiaires,
et l'ATT est ré-estimé sur 200 réplications. La position de l'ATT réel dans
la distribution placebo (centile) et la fraction de pseudo-ATT dépassant un
seuil donné mesurent la crédibilité du résultat. La logique est celle des
tests de permutation appliqués à l'inférence causale
\parencite{heckman1997}.

\subsection{Test d'hétérogénéité des effets}

L'égalité des ATT entre sous-groupes (milieu, sexe, groupe d'âge) est testée
par une interaction triple $t \times D \times G$ dans la régression DD~; la
significativité du coefficient d'interaction indique une différence d'effet
entre modalités du facteur $G$. \textit{Source :} \textcite{card1994}.

\subsection{Analyse de l'attrition différentielle}

La comparaison des ménages suivis et perdus entre les deux vagues repose sur
la différence standardisée moyenne (même formule que ci-dessus) et un test de
Student par covariable, afin de vérifier que l'attrition n'introduit pas de
biais de sélection majeur. \textit{Source :} \textcite{rosenbaum2002}.
